\section{Conclusion}

In many domains, people choose to disseminate information across different modalities and formats for better communication to broader audiences. We proposed a unified view of document transformation and evaluation. We showed that LLMs are capable of templatic document generation with minimal supervision, and that a structured, intermediate representation can improve performance, particularly for smaller models. As LLMs continue to generally improve, we predict that this finding will continue to be true, especially for smaller models. Both summarization and structured document generation are complex tasks, and similar to Chain-of-Thought prompting~\cite{Wei2022ChainOT}, models tend to perform better when complex tasks are broken up into individual steps. This is the core idea behind our methodology. One finding that may not hold to better models is that in our experimentation, we found that allowing LLMs to create their own structure for the intermediate representation do not perform as well. As models improve, they may become better at designing their own intermediate structure, and may even be able to surpass our proposed structure.
We also introduced a flexible precision-recall framework for automatic evaluation that easily integrates existing evaluation metrics into a unified system and allows for comparison across diverse datasets without additional task specific metric design. Finally, we conducted a human evaluation and showed that annotators prefer the documents generated using the intermediate representation 82\% of the time and that our evaluation framework correlates better with human preference than standard evaluation metrics.